%=======================================================[[Define parameter of document]]
\documentclass{article}
\pagestyle{empty}
\setlength{\paperwidth}{8.5in}
\setlength{\paperheight}{11in} 
\setlength{\topmargin}{-0.75in}
\setlength{\headsep}{0.00in} 
\setlength{\headheight}{0.00in}
\setlength{\evensidemargin}{-0.50in}
\setlength{\oddsidemargin}{-0.50in} 
\setlength{\textwidth}{7.50in}
\setlength{\textheight}{10in} 
\setlength{\voffset}{0.00in}
\setlength{\hoffset}{0.00in} 
\setlength{\marginparwidth}{0.00in}
\setlength{\marginparsep}{0.00in} 
\setlength{\parindent}{0.0in}
\setlength{\parskip}{0.0in}

%=======================================================[[Package]]
%\usepackage[dvipdf]{graphicx}
%\usepackage{amsfonts}
%\usepackage{mathptmx}
%\usepackage[scaled=.90]{helvet}
%\usepackage{courier} 
\usepackage{amssymb}
\usepackage{amsmath}
\usepackage{amsthm}
%\usepackage[T1]{fontenc}
\usepackage{linalgjh}
\usepackage{paralist}
\usepackage{verbatim}
\usepackage{hyperref}
\hypersetup{
    colorlinks=true,
    linkcolor=blue,
    filecolor=blue,
    urlcolor=blue,
    citecolor=blue
}
\urlstyle{same}

%\usepackage[margin=0.5in]{geometry}

\usepackage{fancyhdr}
\usepackage{lastpage}
 
\pagestyle{fancy}
\fancyhf{}
\renewcommand{\headrulewidth}{0pt}
\rfoot{Page \thepage \hspace{1pt} of \pageref*{LastPage}}

\def\Dash{\unskip\thinspace\textemdash\thinspace\ignorespaces\allowbreak}

\theoremstyle{plain}
\newtheorem{theorem}{Theorem}[section]
\newtheorem{lemma}[theorem]{Lemma}
\newtheorem{prop}[theorem]{Proposition}
\newtheorem{cor}[theorem]{Corollary}

\theoremstyle{definition}
\newtheorem{definition}{Definition}[section]
\newtheorem{conj}{Conjecture}[section]
\newtheorem{exmp}{Example}[section]

\theoremstyle{remark}
\newtheorem*{remark}{Remark}
\newtheorem*{note}{Note}

\theoremstyle{example}
\newtheorem{example}{Example}[section]

\renewcommand{\vec}[1]{\mathbf{#1}}
%=======================================================[[Title page information]]
%=======================================================[[Document begins here]]
\begin{document}
\begin{center}
\begin{huge}
	Vector Spaces of $\Re^n$
\end{huge}
\end{center}
\tiny
\fbox{\parbox{\textwidth}{\scshape
From \emph{Linear Algebra} by Jim Hefferon 
\texttt{\url{http://joshua.smcvt.edu/linearalgebra}},
{A First Course in Linear Algebra} by Robert A. Beezer
\texttt{\url{http://linear.ups.edu/}}, and \emph{Linear Algebra with Applications (Lyryx)} by Keith W. Nicholson  
modified by Emilie Richer and Yann Lamontagne \texttt{\url{http://obeymath.org}}  (licensed under GNU Free Documentation License or the
Creative Commons Attribution-ShareAlike License.), \today.}}
\normalsize

\stepcounter{section}
\section{Subspaces}

\begin{definition} 
A \definend{subspace} of $\mathbb{R}^n$ is a subset that is closed under addition and scalar multiplication. That is, A set\footnotemark $U$ of vectors in $\Re^n$ is called a \textbf{subspace}\index{subspaces!defined} of $\Re^n$ if it satisfies the following properties:\index{set of all ordered $n$-tuples ($\Re^n$)!subspaces}

\begin{itemize}
\item[S1.] The zero vector $\vec{0} \in U$.\index{vectors!zero vector}\index{zero vector}

\item[S2.] If $\vec{x} \in U$ and $\vec{y} \in U$, then $\vec{x} + \vec{y} \in U$.

\item[S3.] If $\vec{x} \in U$, then $a\vec{x} \in U$ for every real number $a$.

\end{itemize}
\footnotetext{We use the language of sets. Informally, a \textbf{set} $X$ is a collection of objects, called the \textbf{elements} of the set. The fact that $x$ is an element of $X$ is denoted $x \in X$. Two sets $X$ and $Y$ are called equal (written $X = Y$) if they have the same elements. If every element of $X$ is in the set $Y$, we say that $X$ is a \textbf{subset} of $Y$, and write $X \subseteq Y$. Hence $X \subseteq Y$ and $Y \subseteq X$ both hold if and only if $X = Y$.}
\end{definition}

\begin{example}  \label{ex:TrivSbspReFour}
Show that the singleton subset of $\Re^n$
\begin{equation*}
  \{ \vec{0} \}
\end{equation*}
is a subspace.
\end{example}
\vspace{2in}

\begin{definition} \label{df:TrivialVectorSpace}
%<*df:TrivialVectorSpace>
A one-element subspace is a \definend{trivial}\index{vector space!trivial}%
\index{trivial space}
subspace.
%</df:TrivialVectorSpace>
\end{definition}

\begin{example}  \label{cex:RPlusNotSubSp}
Show that the set \( \Re^+ \) is not a subspace of the vector space \( \Re^1 \).
\end{example}
\vspace{1in}

\begin{example}   \label{ex:SubspacesRTwo}
Show that the \( x \)-axis in \( \Re^2 \) is a subspace.
\end{example}
\vspace{1.75in}
\begin{note}
Subspaces of $\Re^2$ are: $\Re^2$, lines through the origin, trivial subspace.
\end{note}
\newpage

\begin{example} 
Show that the subset of \( \Re^3 \) that is a plane through the origin
\begin{equation*}
  P=\left\{ \colvec{x \\ y \\ z}  \Big|\; x+y+z=0\right\}
\end{equation*}
is a subspace.
\end{example}
\vspace{4.2in}
\begin{example}
Determine which subsets are a subspace of $\Re^3$: $W_1=\left\{(x,\;0,\;z)\; |\; x,\;z\in\Re\right\}$, $W_2=\left\{(x,\;1,\;z)\; |\; x,\;z\in\Re\right\}$.
\end{example}
\vspace{4.2in}


\begin{note}
Subspaces of $\Re^3$ are: $\Re^3$, planes through the origin, lines through the origin, trivial subspace.
\end{note}
\newpage
\begin{theorem}
If $A\mathbf{x}=\mathbf{0}$ is a homogenous linear system of $m$ equations in $n$ unknowns, then 
the set of \definend{solution vectors} is a subspace of $\Re^n$.
\end{theorem}
\vspace{4in}

\begin{example}
Express the subspace $W=\left\{\mathbf{x}\;|\;A\mathbf{x}=\mathbf{0}\text{ and }\mathbf{x}\in\Re^3\right\}$ of $\Re^3$ as a parametric equation where
\[
A=\begin{bmatrix}
1 & -1 & 0\\
2 & 1 & 1\\
-1 & 4 & 1
\end{bmatrix}
\]
\end{example}
\newpage
\section{Spanning Sets}
\begin{definition}
A vector $\vec{w}$ is a \definend{linear combination} of a set 
$S = \left\{ \vec{s}_1, \vec{s}_2, \ldots, \vec{s}_n\right\}$ if it 
can be expressed in the form
\[
\vec{w} = c_1\vec{s}_1+c_2\vec{s}_2+\cdots+c_n\vec{s}_n
\]
where $c_i\in\Re$.
\end{definition}

\begin{example}
Any vector $(a, b, c)\in\Re^3$ is expressible as a linear combination of 
the standard basis vectors
\[
(a, b, c) = a\vec{i}+b\vec{j}+c\vec{k}
\]
\end{example}

\begin{example}
Show that $\vec{w}=(3, -4, 2)$ is a linear combination of $\vec{u}=(4, -3, 3)$ and 
$\vec{v}=(1, 1, 1)$
\end{example}
\vspace{2in}

\begin{example}
	Determine whether $\vec{p}=(2,\;-1)$ is a linear combination of $\vec{p}_1=(2,\;1)$ and $\vec{p}_2=(4,\;2)$.
\end{example}
\vspace{2in}

\begin{example}
	Determine if $\vec{x}$ is a linear combination of $\vec{a}$, $\vec{b}$, $\vec{c}$ where
\[
	\vec{x} = (2,\; -1,\;3,\; 1)\;\;\;
	\vec{a} = (2,\; -1,\;3,\; 0)\;\;\;
	\vec{b} = (1,\; 2,\; 3,\; -1)\;\;\;
	\vec{c} = (-4,\; -3,\;-9,\; 3)
\] 
\end{example}
\newpage
\begin{definition} \label{df:Span}
%<*df:Span>
The \definend{span} of a nonempty subset \( S \) of a
vector space is the set of all linear combinations of vectors from \( S \).
\begin{equation*}
  \mbox{span}(S) =\{ c_1\vec{s}_1+\cdots+c_n\vec{s}_n
            \suchthat c_1,\ldots, c_n\in\Re
            \text{\ and\ } \vec{s}_1,\ldots,\vec{s}_n\in S \}
\end{equation*}
The span of the empty subset of a subspace is its trivial subspace.
%</df:Span>
\end{definition}

\begin{theorem}   \label{le:SpanIsASubsp}
%<*lm:SpanIsASubsp>
The span of any subset is a subspace.  In addition, the 
span of a subset is the smallest subspace containing all 
vectors of the subset.
%</lm:SpanIsASubsp>
\end{theorem}
\newpage

\begin{example}
Let $W=\mbox{span}\left\{\vec{u}, \vec{v} \right\}$ where $\vec{u}=(1,\; 0,\; 0)$ 
and $\vec{v}=(0,\; 1,\; 0)$.  Is $(1,\; 1,\; 1)$ in $W$? What about $(2,\; 3,\; 0)$?
\end{example}
\vspace{2.5in}

\begin{example}
Express the subspace $W=\left\{\mathbf{x}\;|\;A\mathbf{x}=\mathbf{0}\text{ and }\mathbf{x}\in\Re^3\right\}$ of $\Re^3$ as a span of a set of vectors where
\[
A=\begin{bmatrix}
1 & -1 & 1\\
2 & -2 & 2\\
-1 & 1 & -1
\end{bmatrix}
\]
\end{example}
\newpage

\begin{example}
Determine whether $S=\left\{(2,\; 1,\; 0),\; (-1,\; 3,\; 1),\; (1,\; 1,\; 1)\right\}$ spans $\Re^3$.
\end{example}
\vspace{5in}

\begin{example}
	Determine whether $S=\left\{(1,\;1,\;1),\; (1,\;0,\;-1)\right\}$ spans $\Re^3$.
\end{example}
\newpage
\begin{example}
	Find the conditions on the coefficient of $\vec{p}=(a_0,\;a_1,\;a_2)$ such that $\vec{p}\in \mbox{span}(\{(1,\;1,\;1),\; (1,\;2,\;3)\})$
\end{example}
\vspace{5in}
\begin{theorem}
Let $S$ and $S'$ be subsets of a subspace $V$. If every vector in $S$ is expressible 
as a linear combination of the vectors in $S'$ then $\mbox{span}(S)$ is a subspace of $\mbox{span}(S')$. 
If in addition every vector of $S'$ is expressible 
as a linear combination of the vectors in $S$ then $\mbox{span}(S)=\mbox{span}(S')$.
\end{theorem}

\begin{example}
	Let $S=\left\{\vec{m}_1, \vec{m}_2 \right\}$ and $S'=\left\{\vec{m}_3, \vec{m}_4\right\}$ where
\[
	\vec{m}_1 = (2,\; -1,\;3,\; 0),\;\;\;
	\vec{m}_2 = (0,\; 1,\;1,\; 0),\;\;\;
	\vec{m}_3 = (4,\; -1,\;7,\; 0),\;\;\;
	\vec{m}_4 = (0,\; 2,\;2,\; 0)
\] 
Is $\mbox{span}(S)=\mbox{span}(S')$?
\end{example}
\newpage
\begin{example}  \label{ex:SubspRThree}
The picture below shows the subspaces of \( \Re^3 \) that we now know of, the 
trivial subspace, the lines through the origin,
the planes through the origin, and the whole space
(of course, the picture shows only a few of the infinitely many subspaces). 
In the next section we will prove that $\Re^3$ has no other
type of subspaces, so in fact this picture shows them all.

That picture describes the subspaces 
as spans of sets with a minimal number of members.
Note that the subspaces 
fall naturally into levels\Dash planes on one level, 
lines on another,
etc.\Dash according to how many vectors are in the minimal-sized
spanning set.
\begin{center}
  \setlength{\unitlength}{4pt}
  \begin{picture}(75,38)(0,-2) %subspaces of R-three
      \thinlines
      \put(45,31){\makebox(0,0)[bl]{
                        \scriptsize \( %\Re^3=
                                   \set{x\colvec[r]{1 \\ 0 \\ 0}
                                              +y\colvec[r]{0 \\ 1 \\ 0}
                                              +z\colvec[r]{0 \\ 0 \\ 1}} \)} }
    %Next the dimension two subspaces
      \put(43,31){\line(-4,-1){25} } % connects R3 to 2,a
      %set 2,a:
      \put(0,22){\makebox(0,0)[l]{\scriptsize\( \set{x\colvec[r]{1 \\ 0 \\ 0}
                                                 +y\colvec[r]{0 \\ 1 \\ 0} }\) }}
      \put(48,29){\line(-3,-1){12} } % connects R3 to 2,b
      %set 2,b:
      \put(20,22){\makebox(0,0)[l]{\scriptsize\( \set{x\colvec[r]{1 \\ 0 \\ 0}
                                                 +z\colvec[r]{0 \\ 0 \\ 1} }\) }}
      \put(53,29){\line(-1,-1){3}} % connects R3 to 2,c
      %set 2,c:
      \put(40,22){\makebox(0,0)[l]{\scriptsize\( \set{x\colvec[r]{1 \\ 1 \\ 0}
                                                 +z\colvec[r]{0 \\ 0 \\ 1} }\) }}
      \put(60,22){\makebox(0,0)[l]{$\cdots$} }
    %Next the dimension one subspaces
      \put(7,18){\line(-1,-4){1} } % connects 2,a to 1,a
      \put(22,18){\line(-3,-1){13} } % connects 2,c to 1,a
      %set 1,a:
      \put(0,10){\makebox(0,0)[l]{\scriptsize\( \set{x\colvec[r]{1 \\ 0 \\ 0}} \)} }
      \put(12,18){\line(1,-2){2} } % connects 2,a to 1,b
      %set 1,b:
      \put(12,10){\makebox(0,0)[l]{\scriptsize\( \set{y\colvec[r]{0 \\ 1 \\ 0}} \)} }
      \put(14,18){\line(2,-1){10} } % connects 2,a to 1,c
      %set 1,c:
      \put(24,10){\makebox(0,0)[l]{\scriptsize\( \set{y\colvec[r]{2 \\ 1 \\ 0}} \)} }
      \put(46,18){\line(-1,-1){3.25} } %connects 2,c to 1,d
      %set 1,d:
      \put(36,10){\makebox(0,0)[l]{\scriptsize\( \set{y\colvec[r]{1 \\ 1 \\ 1}} \)} }
      \put(48,10){\makebox(0,0)[l]{$\cdots$} }
    %Finally, the trivial subspace.
      \put(9,7.5){\line(4,-1){30} } %connects 1,a to trivial
      \put(18.9,7.7){\line(3,-1){20} } %connects 1,b to trivial
      \put(30.4,6.8){\line(2,-1){10} } %connects 1,c to trivial
      \put(41,6){\line(1,-2){1.5} } %connects 1,d to trivial
      \put(45,0){\makebox(0,0){\scriptsize\( \set{\colvec[r]{0 \\ 0 \\ 0} } \)} }
  \end{picture}
\end{center}
\noindent The line segments between levels connect subspaces with 
their superspaces. 
\end{example}
\newpage
%\mbox{ }
%\newpage

\section{Linear Independence}

\begin{definition}\label{def:LinInd}
%<*df:LinInd>
A nonempty subset $S = \left\{ \vec{s}_1,\; \vec{s}_2,\; \ldots,\; \vec{s}_n\right\}$ of a vector space is
\definend{linearly independent}
if 
\[
\vec{0} = c_1\vec{s}_1+c_2\vec{s}_2+\cdots+c_n\vec{s}_n
\]
only has the \definend{trivial solution}, that is, $c_1=c_2=\ldots =c_n=0$. Otherwise it is \definend{linearly dependent}.
%</df:LinInd>
\end{definition}

\begin{example}
Determine if $S=\{\vec{u},\; \vec{v},\; \vec{w}\}$ where
\[
\vec{u}=(2, 1, 3),\;\;\; \vec{v}=(1, 3, 1),\;\;\; \vec{w}=(3, 4, 4)
\] 
is linearly independent.
\end{example}
\vspace{1.5in}

\begin{example}
Determine if $S=\{M_1, M_2, M_3\}$ where
\[
M_1 =
\begin{bmatrix}
1 & 0\\
2 & 0
\end{bmatrix},\;\;\;
M_2 =
\begin{bmatrix}
2 & 1\\
3 & 1
\end{bmatrix},\;\;\;
M_3 =
\begin{bmatrix}
-1 & -3\\
-3 & 4
\end{bmatrix}
\]
is linearly independent.
\end{example}
\newpage

\begin{example}
Suppose $S=\{\vec{v}_1,\; \vec{v}_2,\; \vec{v}_3\}$ is linearly independent.  Are the vectors
\[
\vec{v}_1-\vec{v}_2,\; \vec{v}_3+\vec{v}_2,\; \vec{v}_1-\vec{v}_3
\]
linearly independent?
\end{example}
\vspace{5in}


\begin{theorem}
The polynomials
\[
1,\;x,\;x^2,\;\ldots,\;x^n
\]
form a linearly independent set.
\end{theorem}
\newpage

\begin{example}
Determine if $S=\{p_1(x), p_2(x), p_3(x)\}$ where
\[
p_1(x)=1-2x+x^2,\;\;\; p_2(x)=3+2x^2,\;\;\; p_3(x)=1+x+x^2
\] 
is linearly independent.
\end{example}
\vspace{3in}


\begin{theorem} A subset \( S=\set{\vec{s}_1,\vec{s}_2,\dots,\vec{s}_n} \) of a vector space
is linearly dependent if and only if some \( \vec{s}_i \)
is a linear combination of the vectors $S-\{\vec{s}_i\}$.
\end{theorem}
\newpage
\begin{cor} It follows, that subset \( S=\set{\vec{s}_1,\vec{s}_2,\dots,\vec{s}_n} \) of a vector space
is linearly independent if and only if no \( \vec{s}_i \)
is a linear combination of the vectors $S-\{\vec{s}_i\}$.
\end{cor}


\begin{example} 
What can be said about the linear independence of the following set $S = \left\{ \vec{s}_1,\; \vec{s}_2,\; \ldots,\; \vec{s}_n,\; \vec{0}\right\}$?
\end{example}
\vspace{2in}

\begin{example} 
What can be said about the linear independence of the following set $S = \left\{ \vec{s}_1 \right\}$?
\end{example}
\vspace{2in}

\begin{example} 
What can be said about the linear independence of the following set $S = \left\{ \vec{s}_1,\; \vec{s}_2 \right\}$?
\end{example}
\vspace{2in}


\begin{note}
Geometric Interpretation:\\
Two vectors in $\Re^2$ and $\Re^3$ are linearly independent if and only if the vectors do not lie on the same line
(when their initial points are at the origin).\\

Three vectors in $\Re^3$ are linearly independent if they do not lie on the same plane (when their initial points are at the origin).
\end{note}
\newpage
%\mbox{ }
%\newpage

\section{Coordinates and Basis}

\begin{definition}
A \definend{basis} $B$
for a subspace $V$ is a sequence of vectors $\{\vec{b}_1,\;\vec{b}_2,\;\ldots,\;\vec{b}_n\}$ that is linearly independent 
and that spans the subspace $V$.
\end{definition}
\vspace{1.5in}

\begin{definition}
For $\Re^2$ the \definend{standard} (or \definend{natural}) basis is $\{\vec{i},\; \vec{j}\}$.\\
For $\Re^3$ the \definend{standard} (or \definend{natural}) basis is $\{\vec{i},\; \vec{j},\; \vec{k}\}$.\\
For $\Re^n$ the \definend{standard} (or \definend{natural}) basis is
$\{ \vec{e}_1=(1,\;0,\;0,\;\ldots,\;0),\; \vec{e}_2=(0,\; 1,\; 0,\;\ldots, 0),\; \vec{e}_3=(0,\; 0,\; 1,\;\ldots, 0), \ldots, \vec{e}_n=(0,\; 0,\; 0,\;\ldots, 1)\}$
\end{definition}

\begin{example}
Determine if $B=\{\vec{v}_1,\; \vec{v}_2, \vec{v}_3\}$ where
\[
\vec{v}_1 = (1, 2, 1),\;\;\; \vec{v}_2 = (2, 9, 0),\;\;\; \vec{v}_3=(3, 3, 4)
\]
is a basis for $\Re^3$
\end{example}
\newpage


\begin{theorem}
If $B=\{\vec{b}_1,\; \vec{b}_2,\;\ldots,\; \vec{b}_n\}$ is linearly independent then $B$ is a basis for $V=\mbox{span}(B)$.
\end{theorem}
\vspace{2in}

\begin{theorem}[Uniqueness of Basis Representation] 
In any subspace $V$, a subset $B$ is a basis
if and only if each vector in the subspace $V$ can be expressed as a linear combination of elements of the subset $B$ 
in one and only one way.
\end{theorem}
\vspace{8in}

\begin{definition} 
In a subspace with basis $B$
the \definend{coordinate vector of $\vec{v}$ relative to $B$} (or \definend{representation of \( \vec{v} \) with respect to \( B \)}) is a vector of the coefficients used to express $\vec{v}$ as a 
linear combination of the basis vectors:
\begin{equation*}
  \rep{\vec{v}}{B} = (\vec{v})_B = (c_1,\; c_2,\; \ldots,\; c_n )
\end{equation*}
where
\( B=\{\vec{b}_1,\; \vec{b}_2,\; \dots,\;\vec{b}_n\} \) and 
\( \vec{v}=c_1\vec{b}_1 + c_2\vec{b}_2 + \ldots + c_n\vec{b}_n \).
\end{definition}

\begin{remark}
The definition of the basis requires that a basis be a sequence
because without that we couldn't write these coordinates in a fixed order. 
\end{remark}

\begin{example}
Find the coordinate vector of $\vec{v}=(3,4)$ relative to the basis 
$B=\{\vec{i},\;\vec{j}\}$.
\end{example}
\vspace{1in}

\begin{example}
Find the coordinate vector of $\vec{v}=(3,4)$ relative to the basis 
$B=\{\vec{b}_1,\;\vec{b}_2\}$ where
\[
\vec{b}_1=(2, -4)\;\;\;\vec{b}_2=(3, 2)
\]
\end{example}
\vspace{2in}
\newpage
\section{Dimension}

The previous subsection defines a basis of a vector space and
shows that a space can have many different bases.
So we cannot talk about ``the'' basis for a vector space.
True, some vector spaces have bases that strike us as more natural
than others.  We cannot, in general, associate with a space any single basis that
best describes it.

We can however find something about the bases that
is uniquely associated with the space.
This subsection shows that 
any two bases for a space have the same number of elements.
So with each space we can associate a number, 
the number of vectors in any of its bases.

Before we start, we first
limit our attention to spaces where at least one basis has only finitely
many members.

\begin{definition}
A subspace is \definend{finite-dimensional}
if it has a basis with only finitely many vectors.
\end{definition}

\noindent One space that is not finite-dimensional is $\Re^\infty$
this space is not spanned by any finite subset.

\begin{lemma}[Exchange Lemma] \label{lm:ExchangeLemma}
Assume that 
\( B= \{\vec{b}_1,\; \vec{b}_2,\;\dots,\;\vec{b}_n\} \) is a basis for a
subspace, and that for the vector \( \vec{v} \)
the relationship \( \vec{v}=\lincombo{c}{\vec{b}} \)
has \( c_i\neq 0 \).
Then exchanging \( \vec{b}_i \) for \( \vec{v} \) yields another
basis for the space.
\end{lemma}
\vspace{4.5in}

\begin{theorem}
In any finite-dimensional vector space, all 
bases have the same number of elements.
\end{theorem}

\begin{definition}  \label{df:Dimension}
The \definend{dimension} of a vector space is the number of vectors in any of its bases.
\end{definition}

\begin{example}
Any basis for \( \Re^n \) has \( n \) vectors since the standard basis 
\[
\{ \vec{e}_1=(1,\;0,\;0,\;\ldots,\;0),\; \vec{e}_2=(0,\; 1,\; 0,\;\ldots, 0),\; \vec{e}_3=(0,\; 0,\; 1,\;\ldots, 0), \ldots, \vec{e}_n=(0,\; 0,\; 0,\;\ldots, 1)\}
\]
has \( n \) vectors.
Thus, this definition of `dimension' generalizes the most familiar use of 
term, that $\Re^n$ is $n$-dimensional.
\end{example}

\begin{example}
A trivial space is zero-dimensional since its basis is empty.
\end{example}
\newpage


\begin{example}
Find a basis and determine the dimension of the subspace of $\Re^5$ determined by 
the solution vectors of the following homogeneous system.
\[
\begin{array}{ccccccccccccc}
x_1 & - & 2x_2 & + & x_3 & - & 3x_4 & + & x_5  &  = & 0\\
-x_1 & - & x_2 &  &   & + & x_4  &  &  &  = & 0\\
-2x_1 & + & 3x_2 & + & x_3 &  &   & - & x_5  &  = & 0\\
\end{array}
\]
\end{example}
\vspace{4in}
\begin{cor}
No linearly independent set can have a size greater than the dimension of the
enclosing space.
\end{cor}
\begin{example}
Verify the above corrollary by showing that $S=\{1,\; 1+x,\; 1+x+x^2,\; x+x^2\}$ is
linearly dependent in its enclosing space $P_2$.
\end{example}
\newpage

\begin{cor} 
Any linearly independent set can be expanded to make a basis.
\end{cor}
\begin{example}
Expand the set  $S=\{(1,\;1,\;1,1,1)\}$ to make a basis of the subspace of $\Re^5$ defined by
$W=\{ (a,b,a,b,c)\;|\;a, b, c\in \Re\}$.
\end{example}
\vspace{4in}

\begin{cor} 
Any spanning set can be shrunk to a basis.
\end{cor}
\begin{example}
	Shrink the set $S=\{\vec{m}_1,\; \vec{m}_2,\; \vec{m}_3,\; \vec{m}_4,\; \vec{m}_5\}$ such that it becomes a basis for $\mbox{span}(S)$ where
\[
	\vec{m}_1 = (1,\; 0,\; -2,\; 0),\;\;\;
	\vec{m}_2 = (2,\; 1,\; 3,\; 1),\;\;\;
	\vec{m}_3 = (-1,\; -3,\;4,\; 3),\;\;\;
	\vec{m}_4 = (1,\; 2,\;-5,\; 4),\;\;\;
	\vec{m}_5 = (3,\; -1,\;-6,\; 5)
\]
\end{example}
\newpage


\begin{cor}
In an \( n \)-dimensional space, a set composed of \( n \) vectors is linearly 
independent if and only if it spans the space.
\end{cor}

\begin{example}
	Find a basis for $\Re^4$ containing the linearly independent set $S=\{(1,\;1,\;0,\;0),\;(1,\;0,\;1,0,\;0)\}$ 
\end{example}
\vspace{4in}





\end{document}

